\section{Native code execution and the support of the objectives}
\label{sec:native-loop}
Permitting arbitrary code is not the same as obtaining a working institution,
and obtaining a working institution is not the same as answering an important
optimization question. This section separates these issues in a later native
campaign and a read-only analysis of the earlier completed T/E cohort.
The loss definitions remain unchanged; no participation, novelty or fairness
condition is added to the optimizer's acceptance criteria.

\subsection{A stopped routing-authority campaign}
The new campaign retains the original five fixed Qwen3-8B researchers, the
SciEvo literature and four-year Simulation. It crosses T/E with numerical,
fixed-graph source and code-and-graph design. Two proposals per arm and eight
references are declared in advance. The references include original reviewed
accept-all and typed one- and two-reviewer accept-all institutions, giving
simple controls access to the same new review-routing authority as generated
institutions. These controls change actual review work, not a retrospective
subsampling of already paid reviews.

All eight reference worlds complete. Of twelve declared proposal slots,
eleven are submitted: five complete native execution and six fail technically.
The final numerical T proposal is never submitted because its required fresh
interface calibration fails on the E control. That response is a final HTTP 200
envelope with \code{done=true}, a length termination and 8192 reported output
tokens, but no complete literal submission. The campaign stops at that point;
its unsubmitted slot is neither refilled nor assigned a loss.
Table~\ref{tab:native-stopped-slots} retains all twelve slots. Two cold
stopped-prefix audits agree on the nineteen physical worlds and their original
responses, checkpoints and measurements. This audit scope verifies a stopped
experiment, not normal completion of its planned search and evaluation.

\begin{table}[htbp]\centering\small
\caption{All declared slots in the stopped native campaign. Numbers are
development losses, not unused-seed results. Technical nulls have no imputed
loss; generation labels index each arm, not the global dispatch sequence.}
\label{tab:native-stopped-slots}
\begin{tabular}{@{}llll@{}}
\toprule Objective & Allowed design & Generation 0 & Generation 1\\\midrule
T & Numerical & $-7$ & Not submitted\\
E & Numerical & $-11.167$ & $-11.167$\\
T & Fixed-graph code & $-2$ & Float used as index\\
E & Fixed-graph code & Missing world envelope & $-11.167$\\
T & Code and graph & Missing feature key & Missing world envelope\\
E & Code and graph & Float used as index & Float used as index\\\bottomrule
\end{tabular}\end{table}

Seven of the eight code-permitted submissions change source, and the remaining
one changes numerical router coordinates. None changes connections, node
membership or scheduling. The completed fixed-graph T submission is the
numerical-only edit; the completed fixed-graph E submission changes source.
Therefore neither the method label nor static acceptance establishes that
structural freedom was exercised successfully. The six execution failures
are three float-index errors, two missing output envelopes and one missing
feature-key error. Their literal programs remain unmodified.

The nineteen worlds produce 1796 completed scientist responses, with
1,615,207 reported input and 181,308 output tokens. Separate auxiliary records
contain 256 editor-readiness, 112 staged-proposal, 13 numerical-readiness and
three numerical-proposal intents, plus 25 designer and nineteen scientist
loads. Unreported usage on load calls is unknown, not zero. Readiness and
proposal costs are not silently excluded because the campaign stopped.

\subsection{Separate evaluation of the available submissions}
The stopped campaign has four arms with at least one completed submission;
the two code-and-graph arms have none. A separately frozen study selects the
minimum assigned development loss within each available arm, preserves all
ties, and uses the earliest tied submission. It then evaluates those fixed
policies and six references at seeds 62001, 62011 and 62017, whose prior use is
checked before dispatch. No new designer proposal is generated. Thirty
physical jobs are declared without semantic deduplication. An unavailable arm
has no policy and no score; it does not inherit the seed or a hand reference.

This study evaluates policies available after a stopped prefix. It cannot
complete the missing proposal, establish an equal-budget method comparison or
turn three evaluation seeds into independent design histories. Its results
are integrated only after its own scientific terminal and independent cold
audits close. The new development loop described below does not receive these
unused-seed losses.

\subsection{Execution feedback without host repair}
The execution failures motivate an optional component-testing interface, not
an automatic repair of received programs. A developer bank binds 24 literal
native inputs to three earlier development references, two components and
four saved years. A proposed code artifact can be tested in the existing
network-disabled container; the input, received code, returned value or exact
exception, timeout and charged attempt are retained. Component tests do not
replace native evaluation or define another loss.

Three real-container diagnostics illustrate both utility and limited support.
On the same saved year-one input, the baseline reproduces its recorded output
and a retained missing-envelope program fails. A retained float-index program
returns because that input's paper batch is empty: a passing component test
does not show that the program handles a later nonempty batch. Multi-year
coverage is therefore useful, but still cannot certify a whole trajectory
whose inputs change under a new institution.

A prospective post-edit mode connects literal edits to a fixed suite of four
saved yearly inputs, with at most eight component attempts per proposal.
Errors and compact results enter the next designer prompt; the complete
records remain archived. A failed component test neither vetoes submission
nor causes the host to wrap, revert or replace the generated source. Edits
after the component allowance remain permitted and explicitly untested.
Scripted full-path tests connect calibration, the public certificate API,
the production role factory, the next feedback-bearing request and cold
replay. They show software connectivity, not that the actual designer can
repair its code or improve a native objective.

A new finite development campaign declares four own-history proposals,
three fresh strong references and a seven-world ceiling. Its calibration
declares this optional mode without executing component code. It waits for
the preceding evaluation's independent closure before any model inference.
The reused development seed is not held out. Successful closure, self-repair
and objective improvement remain separate empirical questions.

\subsection{Why different proxies can select the same rule}
Fix the submitted records, their observed review scores and actual review
work $R$. Write $z_i\in\{0,1\}$ for accepting record $i$ and $q_i\geq0$ for
its original mean review score, with $q_i=0$ for an unreviewed record. Then
\begin{equation}
 L_T(z)=-\sum_i z_i,\qquad
 L_E(z)=-\sum_i z_i q_i+\tfrac12 R.
 \label{eq:conditional-accept-all}
\end{equation}
Both conditional losses are minimized by accepting every record; an acceptance
with $q_i=0$ ties under E. Different coefficients therefore need not imply
different acceptance rules. An improvement over the initial threshold rule
alone is weak evidence for a need to invent a complex institution, making
accept-all with the same actual authority an essential control.

This is not a global optimum of the persistent world. An institution can
change reviewer assignments, actual work, future submissions, information and
resource trajectories, which are fixed in the conditional statement.
Moreover, changing the reviewer panel changes the observed $q_i$ itself.
Neither the algebra nor a reduction in E establishes improved scientific
truth. Adding an acceptance penalty, a blind observer or a participation term
would define a new prospective objective, not correct this experiment's loss.

\subsection{Observed annual resource and activity support}
An original-source readout of the earlier closed T/E cohort checks all 176
saved annual checkpoints from its 44 complete physical worlds against the
existing full audit. Five partial failures are retained without complete
outcomes. All 880 observed actor--year boundaries are active; all 44 final
activity stocks are five. Table~\ref{tab:native-resource-support} shows the
individual year-end resource range, not only pooled final totals.

\begin{table}[htbp]\centering\small
\caption{Saved annual support in the earlier 44 complete worlds. Each year
contains 220 actor observations. Extrema can come from different actors and
worlds; rows are not a single actor's trajectory.}
\label{tab:native-resource-support}
\begin{tabular}{@{}lrrr@{}}
\toprule Year & Active / observed & Minimum resources & Maximum resources\\\midrule
1 & 220 / 220 & 90 & 150\\
2 & 220 / 220 & 75 & 200\\
3 & 220 / 220 & 55 & 250\\
4 & 220 / 220 & 30 & 300\\\bottomrule
\end{tabular}\end{table}

Two source-bound readouts and two portable numeric exports agree exactly.
All recorded initial-funding fields are 100; repeated stock fields are not
new funding allocations. The historical within-year resource ledger was
disabled, so saved boundaries cannot exclude temporary inactivity or identify
within-year minima. Nor do 44 completed worlds constitute 44 independent
optimizer histories or show that other institutions could not reduce activity.
The final-activity component alone cannot distinguish these observed completed
worlds. Increasing their seed count is not by itself a remedy for an
unobserved or structurally flat target. A more informative main comparison
must first identify a consequential information, action or timing opportunity
in an extension that retains the persistent researcher world.
\FloatBarrier
